\scanpage{152}

On a
\[
\operatorname{pr}_1^*(F)\otimes\operatorname{pr}_1^*(\unsure{\underline{D}|D}_X F))\unsure{\text{---}|\longrightarrow}\operatorname{pr}_1^*(K_X^\bullet)
\]
\[
\operatorname{pr}_2^*(G)\otimes\operatorname{pr}_{\unsure{2}}^*\unsure{\underline{D}|D}_Y G\unsure{\text{---}|\longrightarrow}\operatorname{pr}_2^*(K_Y)
\]
donc le premier mb de \textcircled{\scriptsize$\alpha$} s'en \unsure{va} dans
\[
\operatorname{pr}_1^*(K_X)\otimes\operatorname{pr}_2^*(K_Y)
\]
il faut l'envoyer \inserted{\unsure{canoniqt}} dans $K_Z^\bullet$. Ceci ne devrait pas offrir de difficulté --- p.ex.\ c'est trivial si $X$ ou $Y$ est lisse/$\unsure{k|S}$.
[en fait, \unsure{on trouve} que le produit tensor est $\sim K_Z^\bullet$].
\margin{ou si $X$ \unsure{et|ou} $Y$ sont lisses sur $\unsure{k}$, \unsure{sinon} \unsure{la|les} \unsure{vérification|vibrations} \unsure{des signes|des cônes} \unsure{inquiétante|inquiétants}\ldots!}
Cela définit bien l'accouplement voulu, d'où un \uline{hom.\ can.}
\margin{\textcircled{**}}
\[
\operatorname{pr}_1^*(F_\bullet)\otimes\operatorname{pr}_2^*(\unsure{\underline{D}|D}_Y(G_\bullet))\longrightarrow\unsure{\underline{D}|D}_Z(\operatorname{pr}_2^*G\otimes\operatorname{pr}_1^*\unsure{\underline{D}|D}_X(F))
\]

\uline{Lemme fondamental de Verdier}: \uline{Cet homomorphisme est un isomorphisme}

\deleted{\uline{Dém.}\ On est ramené comme d'habitude au cas où \inserted{$X,Y$ affines,} $F_\bullet,G_\bullet$ sont de la forme $i_!(A_U)$, $j_!(A_V)$, où $i:U\longrightarrow X$, $j:V\longrightarrow X$ étales, avec $U,V$ affines.}
\margin{151}
